\documentclass[11pt]{article}

\usepackage[T1]{fontenc}
\usepackage{lmodern}
\usepackage[margin=1in]{geometry}
\usepackage{amsmath,amssymb}
\usepackage{booktabs}
\usepackage{graphicx}
\usepackage{xcolor}
\usepackage{microtype}
\usepackage{hyperref}

\hypersetup{
  colorlinks=true,
  linkcolor=blue!45!black,
  citecolor=blue!45!black,
  urlcolor=blue!55!black,
  pdftitle={The gpt5-6 16-component spinor cosmology},
  pdfauthor={Reproducible computational report}
}
\graphicspath{{artifacts/figures/}}
\newcommand{\gptfield}{\ensuremath{\mathrm{gpt5\_6}}}
\newcommand{\dd}{\mathrm{d}}

\title{The \texorpdfstring{\gptfield}{gpt5-6} 16-Component Spinor Cosmology\\
\large Action, stress tensor, CPL reconstruction, numerical solution, and tests}
\author{Reproducible computational report}
\date{17 September 2026}

\begin{document}
\maketitle

\begin{abstract}
We define \gptfield{} as a four-flavor Dirac multiplet in
$\mathbb C^{16}$ with a nonlinear potential depending on the scalar bilinear
$S=\bar\Psi\Psi$.  In a homogeneous flat FLRW spacetime, the Dirac equations
give $S=S_0a^{-3}$ and a perfect-fluid stress tensor with
$\rho_\Psi=U$, $p_\Psi=SU_{,S}-U$, and
$w_\Psi=-1+SU_{,S}/U$.  We reconstruct a potential that realizes the central
CPL history $(w_0,w_a)=(-0.861,-0.60)$ reported for the combined
Unite+BAO+CMB analysis by Camilleri et al.  A DOP853 integration verifies the
bilinear, density, and Friedmann identities at relative errors below
$4.6\times10^{-11}$.  The central reconstructed history crosses $w=-1$ and the
homogeneous null-energy boundary at $a\simeq0.76833$ and begins accelerating at
$z\simeq0.74842$.  We distinguish this result from the separate Unite-only
constant-$w=-0.764$ fit and state the limits of the background reconstruction.
\end{abstract}

\tableofcontents

\section{Scope and claim boundary}

This report supplies a field-theoretic mechanism for a time-varying
dark-energy background and a complete reproducible numerical calculation.  It
does not perform an observational likelihood analysis, propagate parameter
covariances, introduce a new irreducible Lorentz representation, or prove the
stability of perturbations.  The potential is reconstructed from a prescribed
background history.  Consequently, success means that the field equations and
stress tensor consistently realize that history, not that nature selects this
specific microscopic model.

\section{Observational inputs}

The supplied source was a six-page Gmail print discussing the results of
Camilleri et al.\ \cite{camilleri}.  Because it is not the complete paper, the
parameter values and fit labels were checked against the official arXiv v2
record.  The two relevant fits are distinct (Table~\ref{tab:fits}).

\begin{table}[ht]
\centering
\caption{Published constraints used or shown in this project.  Errors are the
reported asymmetric $68.27\%$ credible intervals.}
\label{tab:fits}
\begin{tabular}{@{}llll@{}}
\toprule
Data & Model & $\Omega_m$ & Dark-energy parameters \\
\midrule
Unite only & flat $w$CDM
  & $0.197^{+0.056}_{-0.054}$
  & $w=-0.764^{+0.078}_{-0.096}$ \\
Unite+BAO+CMB & flat $w_0w_a$CDM
  & $0.305\pm0.004$
  & $w_0=-0.861^{+0.044}_{-0.042}$,
    $w_a=-0.60^{+0.17}_{-0.19}$ \\
\bottomrule
\end{tabular}
\end{table}

The Unite-only row assumes a constant equation of state.  The combined row
uses the Chevallier--Polarski--Linder (CPL) form \cite{chevallier,linder}
\begin{equation}
  w(a)=w_0+w_a(1-a).
  \label{eq:cpl}
\end{equation}
The numerical reconstruction uses the central values in the combined row.  The
constant-$w$ result is retained only as a separately labelled comparison.  It
would be statistically incorrect to combine these values as if they arose from
one likelihood.  The source paper also reports weak Bayesian evidence for
evolving dark energy where frequentist summaries reach roughly
$3.1$--$3.3\sigma$; neither statement is evidence specifically for \gptfield{}.

We set $\Omega_{r0}=9\times10^{-5}$ and
$H_0=70\,\mathrm{km\,s^{-1}\,Mpc^{-1}}$ for the demonstration.  These are
numerical choices rather than additional Unite constraints.  The dimensionless
background curves do not depend on the numerical value assigned to $H_0$.

\section{The 16-component field}

An ordinary Dirac spinor in $3+1$ dimensions has four complex components.  We
define
\begin{equation}
  \gptfield\equiv\Psi=
  (\psi_1,\psi_2,\psi_3,\psi_4)^T\in\mathbb C^{16},
  \qquad
  \Gamma^\mu=I_4\otimes\gamma^\mu.
  \label{eq:field}
\end{equation}
For metric signature $(+---)$,
\begin{equation}
  \{\Gamma^\mu,\Gamma^\nu\}=2g^{\mu\nu}I_{16},
  \qquad
  \bar\Psi=\Psi^\dagger\Gamma^0.
\end{equation}
Thus ``16-component'' means a reducible four-flavor multiplet, not a new
irreducible Lorentz spinor.  The flavor-invariant scalar condensate is
\begin{equation}
  S=\bar\Psi\Psi=\sum_{A=1}^{4}\bar\psi_A\psi_A.
\end{equation}

For a homogeneous solution in FLRW, the field equation implies
\begin{equation}
  \dot S+3HS=0,
  \qquad
  S(a)=S_0a^{-3}.
  \label{eq:bilinear}
\end{equation}
The condensate therefore fills space and evolves as the universe expands.  One
positive-$\Gamma^0$ representative can be written
\begin{equation}
  \Psi(t)=a(t)^{-3/2}e^{-i\Theta(t)}\Psi_0,
  \qquad
  \dot\Theta=U_{,S}.
  \label{eq:phase}
\end{equation}
The common phase cancels from $S$, the vector current, and all background
observables used below.  The code function \texttt{gpt5\_6(a)} therefore
returns a convenient phase-suppressed representative with four populated
positive-$\Gamma^0$ components.  It is verified to have 16 components,
$\bar\Psi\Psi=S_0a^{-3}$, and zero spatial current.

\section{Action, field equation, and stress tensor}

Use natural units $c=\hbar=1$, $\kappa=8\pi G$, and a torsion-free spinor
covariant derivative.  The action is
\begin{equation}
  \mathcal S=\int\dd^4x\sqrt{-g}\left[
  \frac{R}{2\kappa}+\mathcal L_\Psi\right]
  +\mathcal S_m+\mathcal S_r,
  \label{eq:action}
\end{equation}
where the Hermitian spinor Lagrangian is
\begin{equation}
  \mathcal L_\Psi=
  \frac{i}{2}\left[
  \bar\Psi\Gamma^\mu\nabla_\mu\Psi
  -(\nabla_\mu\bar\Psi)\Gamma^\mu\Psi
  \right]-U(S).
  \label{eq:lagrangian}
\end{equation}
Independent variation of $\bar\Psi$ gives
\begin{equation}
  i\Gamma^\mu\nabla_\mu\Psi-U_{,S}\Psi=0,
  \label{eq:dirac}
\end{equation}
while variation of $\Psi$ gives the adjoint equation.  Metric or tetrad
variation yields
\begin{align}
T_{\mu\nu}=\frac{i}{4}\big[&
\bar\Psi\Gamma_\mu\nabla_\nu\Psi
+\bar\Psi\Gamma_\nu\nabla_\mu\Psi \nonumber\\
&-(\nabla_\mu\bar\Psi)\Gamma_\nu\Psi
-(\nabla_\nu\bar\Psi)\Gamma_\mu\Psi\big]
-g_{\mu\nu}\mathcal L_\Psi.
\label{eq:stress}
\end{align}
The two field equations imply $\nabla_\mu T^{\mu\nu}=0$ on shell.  This is the
standard nonlinear-spinor construction used in spinor dark-energy models such
as Ref.~\cite{caiwang}, here applied to a four-flavor multiplet.

\section{Kinetic density and thermodynamics}

Define the first-order Dirac kinetic density
\begin{equation}
  K_D=\frac{i}{2}\left[
  \bar\Psi\Gamma^\mu\nabla_\mu\Psi
  -(\nabla_\mu\bar\Psi)\Gamma^\mu\Psi
  \right].
\end{equation}
On the homogeneous equations of motion,
\begin{equation}
  K_D=SU_{,S}.
  \label{eq:kinetic}
\end{equation}
The stress tensor becomes the perfect-fluid expression
\begin{equation}
  T^\mu{}_{\nu}=\operatorname{diag}
  (\rho_\Psi,-p_\Psi,-p_\Psi,-p_\Psi),
\end{equation}
where
\begin{equation}
  \boxed{\rho_\Psi=U(S)},
  \qquad
  \boxed{p_\Psi=SU_{,S}-U(S)},
  \label{eq:rhop}
\end{equation}
and
\begin{equation}
  \boxed{w_\Psi=-1+\frac{SU_{,S}}{U}}.
  \label{eq:w}
\end{equation}
The on-shell Lagrangian equals the pressure,
$\mathcal L_\Psi=p_\Psi$.

The word ``kinetic'' requires care.  $K_D$ is a first-order Dirac density; it is
not a separately nonnegative contribution like $\dot\phi^2/2$ for a canonical
scalar.  In the reconstructed central history $K_D<0$ before the $w=-1$
crossing.  The trajectory therefore must not be called canonical single-field
quintessence.

\section{Exact CPL potential reconstruction}

Conservation of the spinor component gives
\begin{equation}
  \frac{\dd\rho_\Psi}{\dd\ln a}=-3[1+w(a)]\rho_\Psi.
\end{equation}
Inserting Eq.~\eqref{eq:cpl} and integrating produces
\begin{equation}
  \rho_\Psi(a)=\rho_{\Psi0}
  a^{-3(1+w_0+w_a)}\exp[-3w_a(1-a)].
  \label{eq:density}
\end{equation}
Using Eq.~\eqref{eq:bilinear}, $a=(S_0/S)^{1/3}$, defines
\begin{equation}
\boxed{
  U(S)=\rho_{\Psi0}
  \left(\frac{S}{S_0}\right)^{1+w_0+w_a}
  \exp\left[-3w_a\left(
  1-\left(\frac{S_0}{S}\right)^{1/3}\right)\right]
}.
\label{eq:potential}
\end{equation}
Logarithmic differentiation gives
\begin{equation}
  \frac{SU_{,S}}{U}=1+w_0+w_a-w_a
  \left(\frac{S_0}{S}\right)^{1/3}=1+w(a).
\end{equation}
Equations~\eqref{eq:rhop} and \eqref{eq:w} therefore reproduce CPL exactly.
This demonstrates background equivalence, not perturbative equivalence or
uniqueness.

\section{Numerical initial-value problem}

\subsection{Variables and equations}

Let $N=\ln a$ and $\tau=H_0(t-t_0)$.  The state vector is
\begin{equation}
  \mathbf y(N)=\left(\tau,S,\rho_\Psi/\rho_{c0}\right).
\end{equation}
Flatness fixes
\begin{equation}
  \Omega_{\Psi0}=1-\Omega_{m0}-\Omega_{r0}=0.69491.
\end{equation}
The dimensionless Friedmann function and ODEs are
\begin{align}
E^2(N)&\equiv\frac{H^2}{H_0^2}
=\Omega_{r0}e^{-4N}+\Omega_{m0}e^{-3N}
+\frac{\rho_\Psi}{\rho_{c0}},
\label{eq:friedmann}\\
\frac{\dd\tau}{\dd N}&=\frac1E,
&\frac{\dd S}{\dd N}&=-3S,
&\frac{\dd(\rho_\Psi/\rho_{c0})}{\dd N}
&=-3[1+w(e^N)]\frac{\rho_\Psi}{\rho_{c0}}.
\label{eq:odes}
\end{align}

At $N=0$ the initial values are
\begin{equation}
  \tau=0,
  \qquad S=S_0=1,
  \qquad \rho_\Psi/\rho_{c0}=0.69491.
\end{equation}
The code integrates separately from today to $N=-4$ and to $N=1$, then joins
the output in increasing $a$.  The interval
$0.0183\lesssim a\lesssim2.718$ contains 1201 output points.  SciPy's DOP853
method uses relative tolerance $10^{-11}$ and absolute tolerance $10^{-13}$.

\subsection{Discriminating numerical checks}

Three analytic identities independently check the integration:
\begin{equation}
  S=S_0a^{-3},
  \qquad
  \rho_\Psi=U(S),
  \qquad
  E^2=\Omega_{r0}a^{-4}+\Omega_{m0}a^{-3}+U/\rho_{c0}.
\end{equation}
Table~\ref{tab:errors} records the maximum relative discrepancies over all
1201 points.

\begin{table}[ht]
\centering
\caption{Measured numerical consistency errors.}
\label{tab:errors}
\begin{tabular}{@{}lr@{}}
\toprule
Invariant & Maximum relative error \\
\midrule
Bilinear dilution & $4.5232\times10^{-11}$ \\
Density reconstruction & $2.2269\times10^{-11}$ \\
Friedmann closure & $1.7988\times10^{-12}$ \\
\bottomrule
\end{tabular}
\end{table}
These numbers verify the implementation; they are not observational
uncertainties.

\section{Numerical results}

At the present epoch the reconstructed component has
\begin{align}
  \rho_\Psi/\rho_{c0}&=0.69491000,\\
  K_D/\rho_{c0}&=0.09659249,\\
  p_\Psi/\rho_{c0}=\mathcal L_\Psi/\rho_{c0}&=-0.59831751,\\
  w_\Psi&=-0.86100000.
\end{align}
The selected equation-of-state values are
\begin{equation}
  w(0.1)=-1.401,
  \qquad w(1)=-0.861,
  \qquad w(2)=-0.261.
\end{equation}

\subsection{Phantom crossing and the null-energy boundary}

Equation~\eqref{eq:cpl} crosses $w=-1$ at
\begin{equation}
  a_{w=-1}=1+\frac{1+w_0}{w_a}=0.76833333,
  \qquad z=0.30151844.
  \label{eq:crossing}
\end{equation}
Since
\begin{equation}
  \rho_\Psi+p_\Psi=SU_{,S}=\rho_\Psi(1+w),
\end{equation}
this is also the homogeneous null-energy-condition boundary.  The independently
interpolated numerical NEC root is $a=0.76833334$.  The central curve has
$\rho+p<0$ before the root and $\rho+p>0$ after it.  The separately evaluated
constant-$w=-0.764$ benchmark has no $w=-1$ crossing.

\subsection{Acceleration transitions}

The deceleration parameter for the complete radiation, matter, and spinor
background is
\begin{equation}
  q=-\frac{\ddot a}{aH^2}
  =\frac12\left(1+3\frac{p_{\rm tot}}{\rho_{\rm tot}}\right).
\end{equation}
Its roots are
\begin{center}
\begin{tabular}{@{}lrr@{}}
\toprule
Event & $a$ & $z=a^{-1}-1$ \\
\midrule
Acceleration begins & $0.57194538$ & $0.74841869$ \\
Formal future acceleration ends & $1.80120338$ & $-0.44481561$ \\
\bottomrule
\end{tabular}
\end{center}
The first event is in the observed past.  The second follows from extrapolating
the empirical CPL law to $a>1$ and must not be interpreted as an observational
forecast.

\begin{figure}[p]
\centering
\includegraphics[width=0.92\textwidth]{gpt5_6_equation_of_state.pdf}
\caption{The combined-data CPL central reconstruction, the separate
Unite-only constant-$w$ benchmark, and the $\Lambda$CDM line.  Dataset labels in
the legend prevent the two published fits from being conflated.}
\label{fig:w}
\end{figure}

\begin{figure}[p]
\centering
\includegraphics[width=0.92\textwidth]{gpt5_6_energy_budget.pdf}
\caption{Radiation, matter, and reconstructed spinor densities, all normalized
to the present critical density.  Matter is a separate fluid in this run.}
\label{fig:budget}
\end{figure}

\begin{figure}[p]
\centering
\includegraphics[width=0.96\textwidth]{gpt5_6_expansion.pdf}
\caption{Integrated scale factor and total deceleration parameter.  The second
$q=0$ crossing lies in the formal future CPL extrapolation.}
\label{fig:expansion}
\end{figure}

\begin{figure}[p]
\centering
\includegraphics[width=0.92\textwidth]{gpt5_6_potential.pdf}
\caption{The nonlinear potential $U(S)=\rho_\Psi$ and on-shell first-order
Dirac kinetic density $K_D=SU_{,S}$.}
\label{fig:potential}
\end{figure}

\clearpage
\section{Dark energy and dark matter}

\subsection{Dark-energy interpretation}

The spinor condensate is homogeneous in the FLRW approximation, has negative
pressure today, and makes the present universe accelerate when combined with
matter and radiation.  It is therefore a consistent effective dark-energy
mechanism at background level.  Its first-order kinetic structure permits the
central reconstructed history to cross $w=-1$, unlike canonical one-field
quintessence.  This capability does not guarantee viability: a complete
fermionic perturbation analysis is needed to determine whether the NEC-violating
regime is ghost-free, gradient-stable, causal, and weakly coupled.

\subsection{Dark-matter limit}

A different potential demonstrates the mathematical connection to dark matter.
For
\begin{equation}
  U(S)=mS,
\end{equation}
one obtains
\begin{equation}
  p=SU_{,S}-U=0,
  \qquad w=0,
  \qquad \rho=mS\propto a^{-3}.
\end{equation}
The homogeneous component is then dust-like.  A potential
$U=\Lambda+mS$ can mix constant and dust-like terms.  These facts do not make
the CPL reconstruction used here dark matter.  Equation~\eqref{eq:friedmann}
already includes a separate $\Omega_{m0}a^{-3}$ fluid; identifying the same
spinor energy with all matter would double-count it.  A unified dark-sector
model would require a new fit and successful calculations of clustering, sound
speed, CMB anisotropies, lensing, and structure growth.

\section{Limitations and open work}

\begin{itemize}
\item The field is a reducible four-flavor effective multiplet, not a newly
discovered irreducible spin representation.
\item Only a homogeneous scalar condensate is modeled; isotropy is imposed at
the background level.
\item The potential is reverse-engineered from a CPL central curve and is not
unique.
\item No observational likelihood, covariance, nuisance parameters, or
uncertainty bands are calculated.
\item No scalar, vector, tensor, or fermionic perturbations are evolved.
\item Ghost, gradient, causal, radiative, and strong-coupling stability are not
established.
\item Future CPL behavior is an extrapolation outside the fitted domain.
\item $S_0=1$ is a background normalization.  Particle-physics interpretation
requires a dimensionful condensate scale and couplings.
\end{itemize}

\section{Choice of numerical solver}

The requested Linux rustSolveIt repository \cite{rustsolveit} was examined.  It
vendors a pure-Rust translation of SUNDIALS 7.8.0 and includes a
mechanics-oriented notebook language and Jupyter wrapper.  It does not expose a
drop-in Python ODE interface for the custom state in Eq.~\eqref{eq:odes}.
Direct use would require a large external workspace, a separate Rust executable,
and a Python/Rust data bridge.

This system is smooth, nonstiff, has three variables, and has closed-form
invariants.  DOP853 reaches relative consistency errors near $10^{-11}$, so the
extra integration layer offers no numerical benefit for this calculation.  A
future stiff perturbation system could justify a dedicated CVODE implementation.
This is an interface and scope decision, not a claim that SUNDIALS is
unsuitable.

\section{Reproduction guide}

From the repository root on Linux or macOS, create the environment:
\begin{verbatim}
python3 -m venv .venv
.venv/bin/python -m pip install --upgrade pip
.venv/bin/python -m pip install -r requirements.txt
\end{verbatim}
Run the focused tests and regenerate numerical artifacts:
\begin{verbatim}
.venv/bin/python -m pytest tests/test_gpt5_6_cosmology.py -q
.venv/bin/python scripts/run_cosmology.py
\end{verbatim}
Build and execute the notebook, then export HTML:
\begin{verbatim}
.venv/bin/python scripts/build_cosmology_notebook.py
.venv/bin/python -m jupyter nbconvert --to notebook --execute \
  --inplace notebooks/gpt5_6_cosmology.ipynb \
  --ExecutePreprocessor.timeout=180 \
  --ExecutePreprocessor.record_timing=False
.venv/bin/python -m jupyter nbconvert --to html \
  --output-dir build notebooks/gpt5_6_cosmology.ipynb
\end{verbatim}
Build this PDF:
\begin{verbatim}
bash scripts/build_documentation.sh
\end{verbatim}
The expected focused result is 10 passing tests, no notebook error outputs,
three invariant errors below $2\times10^{-9}$, four PNG/PDF figure pairs, and
the compiled file \texttt{docs/gpt5\_6\_cosmology.pdf}.  The exact environment,
source hashes, and validation commands are recorded in
\texttt{PROVENANCE.md}.

\begin{thebibliography}{9}
\raggedright

\bibitem{camilleri}
R. Camilleri et al.,
\emph{Supernovae Unite: Combining Pantheon+ and DES-SN5YR},
arXiv:2609.05053v2, especially Sections 7.0.3--7.0.4,
\url{https://arxiv.org/abs/2609.05053}.

\bibitem{caiwang}
Y.-F. Cai and J. Wang,
\emph{Dark Energy Model with Spinor Matter and Its Quintom Scenario},
Class. Quantum Grav. \textbf{25} (2008) 165014,
arXiv:0806.3890,
\url{https://arxiv.org/abs/0806.3890}.

\bibitem{chevallier}
M. Chevallier and D. Polarski,
\emph{Accelerating Universes with Scaling Dark Matter},
Int. J. Mod. Phys. D \textbf{10} (2001) 213--224.

\bibitem{linder}
E. V. Linder,
\emph{Exploring the Expansion History of the Universe},
Phys. Rev. D \textbf{68} (2003) 083504,
arXiv:astro-ph/0208512.

\bibitem{rustsolveit}
once-ere,
\emph{rustSolveIt on pure-Rust SUNDIALS 7.8.0 for Ubuntu Linux},
\url{https://github.com/once-ere/rustSolveIt_linux_SUNDIALS_7_8_0},
accessed 17 September 2026.

\end{thebibliography}

\end{document}
